\section{Complete spectrum discretization and step-function collapse}\label{sec:spectrum}

\begin{corollary}[Initial spectral gap]\label{thm:spectral}
For every $x\in\R$, either $\Lam(x)=0$ or $\Lam(x)\ge1/4$. Thus
$\Lam(\R)\subseteq\{0\}\cup[1/4,\infty]$, and $\Lam(x)=0\iff x\in L$.
\end{corollary}

This is an immediate consequence of Theorems~\ref{thm:class}
and~\ref{thm:gap}; we state it separately because it is the form used
below.

\begin{proof}
If $x\in L$, Theorem~\ref{thm:class} gives $j\|Q_jx\|\to0$, hence $\Lam(x)=0$.
If $x\notin L$, Theorem~\ref{thm:gap} gives $\Lam(x)\ge1/4$. The dichotomy and
the characterization of the zero value follow.
\end{proof}

\begin{remark}\label{rem:linear-indep}
In particular $\Lam(\ee)=0$, since $\ee=2\cdot\frac{\ee+1}{2}-1\in L$; $\ee$ is
\emph{not} a generic point for this exponent. By Corollary~\ref{thm:spectral}, the statement
``$\Lam(B\pi)>0$ for every integer $B\ge1$'' is equivalent to
$B\pi\notin L$ for all $B\ge1$, i.e.\ to the linear independence of
$1,\ee,\pi$ over $\Q$, which is open. Nothing in this article decides it.
\end{remark}

We now show that the initial gap $(0,1/4)$ is merely the bottom rung of an
infinite ladder of equidistant spectral gaps: the entire finite spectrum of
$\Lam$ is quantized in multiples of $1/4$.

\begin{theorem}[Complete discretization of the spectrum]\label{thm:full-spectrum}
For every $x\in\R$ with $\Lam(x)<\infty$, one has $u_j=O(1/j)\to0$ and
\begin{equation}\label{eq:digit-asymp}
\lim_{j\to\infty}\bigl(d_j-4ju_j\bigr)=0,\qquad
\lim_{j\to\infty}\bigl(|d_j|-4j\|Q_jx\|\bigr)=0.
\end{equation}
Consequently $4\Lam(x)=\limsup_{j\to\infty}|d_j|\in\Z_{\ge0}$ whenever
$\Lam(x)<\infty$, and the full spectrum of $\Lam$ is given by the exact identity
\begin{equation}\label{eq:full-spectrum}
\Lam(\R)=\frac14\Z_{\ge0}\cup\{\infty\}
=\left\{0,\frac14,\frac12,\frac34,1,\dots\right\}\cup\{\infty\}.
\end{equation}
In particular, for every integer $k\ge0$, the open interval
$\bigl(\frac{k}{4},\frac{k+1}{4}\bigr)$ is an unconditional empty spectral gap
for $\Lam$.
\end{theorem}

\begin{proof}
If $\Lam(x)<\infty$, then $\|Q_jx\|=|u_j|\le(\Lam(x)+1)/j$ for all large $j$,
so $u_j\to0$. Rearranging \eqref{eq:defect-id} gives
$d_j-4ju_j=2u_j+u_{j-1}-u_{j+1}\to0$, and the reverse triangle inequality
$\bigl||d_j|-4j|u_j|\bigr|\le|d_j-4ju_j|$ yields \eqref{eq:digit-asymp}.
Since $|d_j|\in\Z_{\ge0}$, the sequence $(|d_j|)$ can only have a nonnegative
integer or $\infty$ as its limit superior. By \eqref{eq:digit-asymp},
$4\Lam(x)=\limsup_{j\to\infty}4j\|Q_jx\|=\limsup_{j\to\infty}|d_j|\in\Z_{\ge0}\cup\{\infty\}$,
which excludes every open interval $(k/4,(k+1)/4)$ from $\Lam(\R)$.
Conversely, $\Lam(0)=0$, $\Lam=\infty$ holds on a set of full Lebesgue measure
(Section~\ref{sec:finite-levels}), and for each integer $k\ge1$ choosing a
sufficiently large starting index $N$ and constant digits $d_j\equiv k$ for
$j\ge N$ (as in the proof of Theorem~\ref{thm:gauge} below) produces a point
$x_k\in\R$ with $|u_j|<1/8$ for $j\ge N$, $M_j=\mathrm{round}(Q_jx_k)$, and
$4j\|Q_jx_k\|\to k$, so $\Lam(x_k)=k/4$. Thus every value in
$\frac14\Z_{\ge0}\cup\{\infty\}$ is attained.
\end{proof}

For $c\ge0$ write $E_c:=\{x\in\R:\Lam(x)\le c\}$.

\begin{corollary}[Sub-threshold level sets and gap constancy]\label{thm:sub}
For every integer $k\ge0$ and every $c\in\bigl[\frac{k}{4},\frac{k+1}{4}\bigr)$,
one has $E_c=E_{k/4}=\{x\in\R:u_j(x)\to0\text{ and }\limsup_{j\to\infty}|d_j(x)|\le k\}$.
In particular, if $c<1/4$ ($k=0$), then $E_c=E_0=L$, so $E_c$ is countable and
$\dimH E_c=0$.
\end{corollary}

\begin{proof}
Clearly $E_{k/4}\subseteq E_c$ for $c\ge k/4$. Conversely, if $x\in E_c$ with
$c<(k+1)/4$, then $\Lam(x)\le c<(k+1)/4$ cannot lie in the empty spectral gap
$(k/4,(k+1)/4)$ by Theorem~\ref{thm:full-spectrum}, so $\Lam(x)\le k/4$ and
eventually $|d_j(x)|\le k$. For $k=0$, $\Lam(x)=0\iff x\in L$ by
Corollary~\ref{thm:spectral}, and $L$ is a countable image of $\Z^2$, hence has
Hausdorff dimension $0$ \cite{Falconer}.
\end{proof}

